\par\noindent\begin{minipage}{\linewidth}
\small
\setlength{\tabcolsep}{4pt}
\renewcommand{\arraystretch}{1.13}
\captionof{table}{Comprobación del significado de las medidas en las nueve nubes sintéticas.}\label{tab:S13}
\begin{tabular}{@{}>{\raggedright\arraybackslash}p{4.7cm}>{\raggedright\arraybackslash}p{2.25cm}>{\raggedright\arraybackslash}p{2.55cm}>{\raggedright\arraybackslash}p{2.7cm}>{\raggedright\arraybackslash}p{2.5cm}@{}}
\toprule
\textbf{Nube} & \textbf{Brecha k25} & \textbf{Radio 90/50} & \textbf{Arista máx. / mediana} & \textbf{Componente mayor mutuo} \\
\midrule
Nube redonda & 0,499 & 1,121 & 1,371 & 96,6\% \\
Nube estrecha & 0,503 & 1,132 & 1,400 & 95,5\% \\
Nube ancha & 0,487 & 1,084 & 1,262 & 99,8\% \\
Nube alargada & 0,091 & 1,115 & 1,299 & 99,1\% \\
Nube de rango bajo & 0,098 & 1,589 & 4,913 & 99,9\% \\
Dos grupos separados & 0,000 & 1,766 & 2,253 & 49,5\% \\
Cuatro grupos separados & 0,000 & 1,003 & 1,587 & 24,9\% \\
Dos grupos con puentes & 0,007 & 1,005 & 1,428 & 97,8\% \\
Nube con puntos atípicos & 0,499 & 1,123 & 4,770 & 95,7\% \\
\bottomrule
\end{tabular}
\end{minipage}\par
\par\smallskip{\small Cada nube guardada contiene 1.000 puntos. Una brecha menor puede aparecer sin grupos separados; los cocientes de radios y aristas pueden responder a puntos atípicos o nubes alargadas. El grafo de unión k25 está conectado en los 260 casos principales de modelo y área. La conexión describe el grafo y no queda validada como fragmentación temática general. Las cinco selecciones de igual tamaño conservan 51 alertas de conexión (Tabla S12).\par}
\par\smallskip{\small Se incluyen las 780 filas de referencias gaussianas. Sus errores respecto a lo observado van de -4,69\% a 4,60\% en apertura y de -10,47\% a 19,78\% en PR. Los muestreos gaussianos conservan los momentos ajustados solo en esperanza antes de renormalizar; no son pruebas nulas calibradas. Se adjuntan controles n/k completos; el de fracción fija de vecinos se añadió durante la ejecución.\par}
